% framed knot への refined large N 双対の拡張と，トーラス結び目を超えた正値性の検証（日本語の準備稿）
% 構成：docs/paper/outline.md，主張の格付け：docs/paper/claims.md（各主張の直前に % [ID] で対応を書く），記号：docs/paper/notation.md
% コンパイル：cd docs/paper/ja && latexmk -lualatex main.tex
\documentclass[11pt,a4paper]{ltjsarticle}
\usepackage{amsmath,amssymb,amsthm}
\usepackage{booktabs,array,longtable,graphicx}
\usepackage[margin=25mm]{geometry}
\usepackage[hidelinks]{hyperref}

\newtheorem{theorem}{定理}[section]
\newtheorem{proposition}[theorem]{命題}
\newtheorem{lemma}[theorem]{補題}
\newtheorem{corollary}[theorem]{系}
\newtheorem{conjecture}[theorem]{予想}
\theoremstyle{definition}
\newtheorem{definition}[theorem]{定義}
\newtheorem{hypothesis}[theorem]{仮説}
\newtheorem{compresult}[theorem]{計算結果}
\theoremstyle{remark}
\newtheorem{remark}[theorem]{注意}

\newcommand{\rCS}{\mathrm{rCS}}
\newcommand{\ba}{\mathbf{a}}
\newcommand{\bq}{\mathbf{q}}
\newcommand{\bt}{\mathbf{t}}
\newcommand{\hN}{\widehat{N}}
\newcommand{\hf}{\widehat{f}}
\newcommand{\hbN}{\widehat{\mathbf{N}}}
\newcommand{\fr}{\mathrm{fr}}
\newcommand{\Exp}{\mathrm{Exp}}

\title{framed knot への refined large $N$ 双対の拡張と\\トーラス結び目を超えた正値性の検証\\
\large（日本語の準備稿）}
\author{亀山 昌也（Masaya Kameyama）}
\date{2026年10月9日（準備稿）}

\begin{document}
\maketitle
\begin{abstract}
Kameyama--Nawata~\cite{Kameyama17}（以下 KN）は，refined Chern--Simons 不変量を resolved conifold 上の refined BPS 不変量に書き換え，
それらが非負整数になるという正値性予想を，framing 0 のトーラス結び目・トーラス絡み目と，いくつかの thin な結び目で確かめた．
本論文では，この枠組みを framing をもつ結び目に拡張し，KN が用いたデータを二つの方向に超えて，
$G$ 展開の前の係数 $\hbN$ と後の refined BPS 不変量 $\hN$ の正整数性を検証する．

まず，framing $x$ の不変量を $\rCS^{(x)}_\lambda=\fr[\lambda]^x\rCS_\lambda$ で定める．
$\fr[\lambda]$ は Taki の refined topological vertex の framing の因子（の逆数）に一致し，$t=q$ で Mariño--Vafa の framing の規則に戻る．
$\bt^{-|\rho|x}\hf^{(x)}_\rho$ は $x$ の指数和であり，KN の枠組みの構造的な性質（転置対称性，Laurent 性，$G$ 展開が閉じる条件）は $x$ によらない．
Mariño--Vafa の符号 $(-1)^{|\lambda|x}$ は，Laurent 性を保つために必要であることを示す．

次に，結び目と framing の両方の方向で検証する．
(i) ホモロジーを入力とする限り，任意の結び目・任意の framing で基本表現 $\rho=\yng(1)$ の正値性が成り立つ（命題）．11 交点以下の全素結び目 801 個で確かめた．
(ii) 2 箱では，Laurent 性の必要条件 $\dim H^{\yng(2)}=(\dim H^{\yng(1)})^2$ を導き，予想式を入力として，ツイスト結び目の全ての撚りの数での正値性を示す．
また 8 個の結び目では，全ての整数の framing で正値性を示す．
(iii) 3 箱では，トーラス結び目 7 個と代数的な iterated torus 結び目 7 個（3 本ケーブルと 2 段の反復を含む）について，
定義どおりの refined 不変量から，$|\rho|\le 3$ の全ての $\rho$ で $G$ 展開の前後とも正であることを計算で確かめた．
最大の例は約 183 万項である．
これらに基づき，framed knot の refined LMOV 不変量の正値性を予想として提出する．
\end{abstract}

\tableofcontents

\section{序論}\label{sec:intro}

\subsection{refined large $N$ 双対と KN の正値性予想}
$U(N)$ Chern--Simons 理論の large $N$ 極限は resolved conifold $X=\mathcal{O}(-1)\oplus\mathcal{O}(-1)\to\mathbb{P}^1$ 上の位相的弦理論と等価であり，
結び目 $K$ を加えると，彩色 HOMFLY-PT 多項式は $X$ の中の Lagrangian $L_K$ に端をもつ BPS 状態の数（LMOV 不変量）で書ける
\cite{Ooguri99,Labastida01}．
Kameyama--Nawata~\cite{Kameyama17} は，この描像を refined Chern--Simons 理論~\cite{Aganagic11} に広げた．
M 理論の refined index と Schwinger の計算から自由エネルギーの形を定め，refined CS 不変量 $\rCS_\lambda(K;a,q,t)$ から
refined reformulated invariants $f^q_\mu$, $f^{\bar t}_\mu$ と $\hf_\rho$ を作る式を与えた．
そのうえで，$\hf_\rho$ を太字の変数 $(\ba,\bq,\bt)$ で書いたときの係数 $\hbN_{\rho;i,j,k}$（$G$ 展開の前）と，
$G=(\bq\bt-\bq^{-1}\bt^{-1})(\bq-\bq^{-1})$ の冪に展開した後の refined BPS 不変量 $\hN_{\rho,g,\beta,J_r}$ が，非負整数になることを予想した．
物理的には，この正値性は追加の $U(1)_R$ 対称性が BPS 状態に自明に作用することから期待される．

KN がこの予想を確かめたのは，framing 0 の次のデータである（以下「KN のデータ」と呼ぶ）．
トーラス結び目・トーラス絡み目の $|\lambda|\le 3$ の彩色，
$\mathbf{4_1},\mathbf{5_2},\mathbf{6_1},\mathbf{6_3}$ などの thin な結び目の $\yng(1),\yng(2),\yng(1,1)$（彩色 HOMFLY-PT ホモロジーの Poincaré 多項式を入力とする），
および KN が構成した $\mathbf{4_1}$ の $\yng(2,1)$ の候補である．
本論文では，KN のデータの再現は規約を定めるためだけに用い，成果としては扱わない．

\subsection{本論文の内容}
本論文の目的は，KN の枠組みを framing をもつ結び目に拡張し，KN のデータを二つの方向に超えて正整数性を検証することである．
一つは結び目の方向で，トーラス結び目と thin な結び目から，代数的な iterated torus 結び目，結び目表の全ての結び目，ツイスト結び目・pretzel 結び目の族，thick な結び目へ広げる．
もう一つは framing の方向で，$x=0$ から，有限の範囲の $x$，さらに全ての整数 $x$ へ広げる．

\paragraph{framed knot への拡張（\S\ref{sec:framing}）．}
framing $x$ の不変量を $\rCS^{(x)}_\lambda=\fr[\lambda]^x\,\rCS_\lambda$，
\[
\fr[\lambda]=\bigl(-(q/t)^{1/2}\bigr)^{|\lambda|}\,q^{n(\lambda')}\,t^{-n(\lambda)}
\]
で定める．この因子は Taki~\cite{Taki07} の refined topological vertex の framing の因子の逆数にちょうど一致し，
$t=q$ では Mariño--Vafa~\cite{Marino01} の framing の規則 $W_R\mapsto(-1)^{\ell p}q^{p\kappa_R/2}W_R$ に $x=p$ で戻る．
太字の変数では $\fr[\lambda]$ は符号のない単項式 $\bt^{|\lambda|+2n(\lambda')}\bq^{2(n(\lambda')-n(\lambda))}$ になり，
framing はホモロジーの次数のずらしとして働く．
主な構造的な結果は次のとおりである．
\begin{itemize}
\item $\bt^{-|\rho|x}\hf^{(x)}_\rho$ は $x$ の指数和である（命題~\ref{prop:expsum}）．
\item 転置対称性，Laurent 性，$G$ 展開が閉じる条件は $x$ によらない（命題~\ref{prop:x-indep}）．
\item Mariño--Vafa の符号 $(-1)^{|\lambda|x}$ を外すと，奇数の $x$ で 2 箱の $\hf$ が Laurent でなくなる（命題~\ref{prop:sign}）．
\item KN の単項式の正規化のうち彩色によって変わる部分は，ちょうど Seifert framing の分の framing である（命題~\ref{prop:N2-framing}）．
\end{itemize}

\paragraph{検証（\S\ref{sec:one}--\S\ref{sec:three}）．}
入力の確からしさは結び目のクラスによって異なる（表~\ref{tab:overview}）．
\begin{itemize}
\item 基本表現 $\rho=\yng(1)$：入力が非負係数で DGR の対称性をもてば，$G$ 展開は整数で閉じ，前後とも非負である（命題~\ref{prop:one}）．
DGR の対称性は Gorsky--Hogancamp--Mellit~\cite{Gorsky21} の定理なので，ホモロジーを入力とする限り，任意の結び目・任意の framing で $\yng(1)$ は正である．
\item 2 箱：$\hf_{\yng(2)}$ が Laurent であるためには $\dim H^{\yng(2)}=(\dim H^{\yng(1)})^2$ が必要である（命題~\ref{prop:two-necessary}）．
GNSSS~\cite{Gukov15} の予想式を入力として，ツイスト結び目の全ての撚りの数で 2 箱が正であることを示す（定理~\ref{thm:twist}）．
また，8 個の結び目では全ての整数の framing で 2 箱が正である（定理~\ref{thm:two-all-x}）．
\item 3 箱：トーラス結び目 7 個と，代数的な iterated torus 結び目 7 個で，$|\rho|\le3$ の全ての $\rho$ について $G$ 展開の前後とも正である（計算結果~\ref{cr:three}）．
後者は Cherednik--Danilenko~\cite{Cherednik14} の DAHA による定義どおりの refined 不変量を入力とし，
3 本ケーブル $K7$ と 2 段の反復 $K10$（3 箱で約 183 万項）を含む．
入力の $\rCS_{\yng(2,1)}$ 自体は負の係数をもつが，BPS の段では正になる．
\end{itemize}
これらに基づき，framed knot の refined LMOV 不変量の正値性を予想~\ref{conj:A} として，
代数的な iterated torus 結び目の正値性を予想~\ref{conj:B} として提出する．
双曲結び目の 3 箱には，定義に基づく $\yng(2,1)$ の refined 不変量が無い．構造的な条件では $\yng(2,1)$ が一意に決まらないことを \S\ref{sec:21} で述べる．

\begin{table}[t]
\centering
\caption{検証の全体像．○ は $G$ 展開の前後とも正，— は未実施，× は入力が無い．}\label{tab:overview}
\resizebox{\textwidth}{!}{%
\begin{tabular}{llccccc}
\toprule
結び目 & 入力 & $\yng(1)$ & 2 箱 $x=0$ & 2 箱 $x\neq0$ & 3 箱 $x=0$ & 3 箱 $x\neq0$ \\
\midrule
自明な結び目 & -- & 全 $x$ & ○ & ○ & ○ & $|x|\le6$ \\
$T_{2,3},T_{2,5},T_{2,7},T_{3,4}$ & EHA & 全 $x$ & ○ & ○ & ○ & $|x|\le 3$--$6$ \\
$T_{3,5},T_{2,11},T_{3,7},T_{4,5}$ & EHA & 全 $x$ & ○ & $|x|\le10$ & ○ & $|x|\le10$ \\
$K1,K3$ & DAHA & 全 $x$ & ○ & ○ & ○ & $|x|\le2$（$K1$ は $x=\pm6$ も） \\
$K6a,K7,K8,K9$ & DAHA & 全 $x$ & ○ & $|x|\le10$ & ○ & $|x|\le10$ \\
$K10$ & DAHA & 全 $x$ & ○ & $|x|\le10$ & ○ & $|x|\le2$ \\
$\mathbf{4_1}$ & ホモロジー（$\yng(2,1)$ は KN の候補） & 全 $x$ & ○ & ○ & ○ & $|x|\le4$ \\
結び目表 801 個 & ホモロジー & 全 $x$ & × & × & × & × \\
ツイスト $K_p$ & GNSSS の予想式 & 全 $x$ & 全 $p$ & 注 1 & × & × \\
pretzel $K_{s,t}$ & GNSSS の予想式 & 全 $x$ & 注 2 & -- & × & × \\
$\mathbf{10_{139}},\mathbf{10_{152}}$ & quiver ＋ KRSS の予想式 & 全 $x$ & ○ & -- & × & × \\
\bottomrule
\end{tabular}}

\smallskip
\raggedright
注 1：$p$ と $x$ が同符号の向きでは全て正．逆符号は $|p|\le6$，$|x|\le14$ の直接計算で正．\\
注 2：同符号は全て正．異符号は $\min(|s|,t)\le20$ で証明，$|s|\le30$ かつ $t\le40$ で確認．
\end{table}

\subsection{関連する研究}
非 refined の framed knot の LMOV 型の整数性は Mariño--Vafa~\cite{Marino01} に始まり，
Zhu らの一連の研究~\cite{Zhu21,Luo16,Luo21,Wang25,Zhu25} がある．
Zhu~\cite{Zhu21} の ``refined LMOV conjecture'' の refinement は KN の refined（$t\neq q$）とは別の意味であり（同論文の Remark），本論文の主張とは対象が異なる．
refined の open BPS 不変量の非負性は，単一のブレーン・対称表現について Cheng--Sułkowski~\cite{Cheng21}（framing は固定）と
Banerjee--Ishtiaque--Jeong~\cite{Banerjee26} が論じている．
knots--quivers 対応~\cite{Kucharski17,Kucharski17b,Ekholm18} では，framing の変化は quiver 行列のずらし $C\mapsto C+f$ として現れ，
対称表現の refined LMOV 不変量の整数性が quiver の構造から従う．
本論文の調べた範囲では，KN の $\hf_\rho$ を framing をつけて作り，非対称な $\rho$（$\yng(2,1)$ を含む）まで $G$ 展開の前後の正整数性を調べた研究は見当たらない．
ケーブルの refined 不変量の定義は Cherednik--Danilenko~\cite{Cherednik14,Cherednik15} による．
その明示例は $K1$ の $\yng(1),\yng(2)$ などに限られ，$|\lambda|=3$ の値と refined BPS 不変量は文献に見当たらない．

\subsection{構成}
\S\ref{sec:prelim} で変数と KN の refined LMOV 不変量を復習する．
\S\ref{sec:framing} で framed knot への拡張と，$x$ によらない構造を述べる．
\S\ref{sec:inputs} で入力（トーラス結び目，ケーブル，ホモロジー，予想式）と計算法を説明する．
\S\ref{sec:one}，\S\ref{sec:two}，\S\ref{sec:three} で，それぞれ $\yng(1)$，2 箱，3 箱の検証を述べる．
\S\ref{sec:discussion} で未解決の問題を議論する．
付録に計算法の詳細，非代数的なケーブル，文献の誤植，データとコードをまとめる．

主張の確からしさを区別するため，解析的に導いたものを定理・命題・補題とし，計算機で確かめたものは「計算結果」として範囲とともに書く．
予想式を入力とする主張は，その旨を文中に明記する．

\section{準備}\label{sec:prelim}

\subsection{変数}
KN~\cite{Kameyama17} に従い，$q=e^{i\epsilon_1}$，$t=e^{i\epsilon_2}$ を $\Omega$ 背景 $(z_1,z_2)\mapsto(qz_1,t^{-1}z_2)$ の変数とし，
$SU(N)$ の不変量の large $N$ 極限を $a=t^N$ で表す．
太字の変数 $(\ba,\bq,\bt)$ を
\begin{equation}\label{eq:cov}
a=-\ba^2\bt,\qquad q^{1/2}=-\bq\bt,\qquad t^{1/2}=\bq
\end{equation}
で定める（KN の式 (cov)）．これは Dunfield--Gukov--Rasmussen~\cite{Dunfield05} の superpolynomial の変数と同じであり，
$t=q$ は $\bt=-1$ に当たる．Aganagic--Shakirov~\cite{Aganagic11} の太字の変数も同じである．
Cherednik--Danilenko~\cite{Cherednik14} の DAHA の変数 $(q,t,a)$ は KN の $(q,t)$ と同じで，$a_{\mathrm{DAHA}}=-a$ である．
半整数冪の分枝は \eqref{eq:cov} で固定する．$\hf_\rho$ の全体の符号は規約であり，最低次の項が正になる向きにとる．

\subsection{Macdonald 多項式と自明な結び目}
$P_\lambda(x;q,t)$ を monic な Macdonald 多項式とし，内積 $\langle p_\lambda,p_\mu\rangle_{q,t}=\delta_{\lambda\mu}z_\lambda\prod_i(1-q^{\lambda_i})/(1-t^{\lambda_i})$ に関する
$g_\lambda=(q/t)^{|\lambda|}/\langle P_\lambda,P_\lambda\rangle_{q,t}$ を用いる．
自明な結び目の不変量は主特殊化
\[
P_\lambda(1,t,\dots,t^{N-1})=\prod_{s\in\lambda}\frac{t^{l'(s)}-a\,q^{a'(s)}}{1-q^{a(s)}t^{l(s)+1}}\qquad(a=t^N)
\]
を対称化したものである．$a(s),l(s)$ は腕と脚の長さ，$a'(s),l'(s)$ は余腕と余脚の長さである．
$n(\lambda)=\sum_i(i-1)\lambda_i$ とおき，
\[
T_\lambda=q^{n(\lambda')}\,t^{-n(\lambda)}
\]
を refined の framing の固有値と呼ぶ．$t=q$ では $q^{\sum_{s\in\lambda}c(s)}$（$c(s)$ は content）になる．

\subsection{refined CS 不変量}
KN と同じく，unreduced の refined CS 不変量を $\overline{\rCS}_\lambda(K;a,q,t)$，reduced の不変量を $\rCS_\lambda(K;a,q,t)$ と書く．両者は
\[
\overline{\rCS}_\lambda(K;a,q,t)=\overline{\rCS}_\lambda(\bigcirc;a,q,t)\,\rCS_\lambda(K;a,q,t)
\]
で結ばれる．ここで $\bigcirc$ は自明な結び目である．
トーラス結び目では，Aganagic--Shakirov の $S,T$ 行列による定義~\cite{Aganagic11}，Cherednik の DAHA による定義~\cite{Cherednik11}，
Gorsky--Neguţ~\cite{Gorsky13b} の楕円 Hall 代数による表示が知られている．
本論文では後者で計算する（\S\ref{sec:inputs}）．
これらの定義は単項式倍を除いてしか値を決めない．KN の BPS 抽出は単項式倍で結果が変わるので，単項式の正規化を \S\ref{sec:normalization} で固定する．
$\U(N)$ と $\SU(N)$ の違いは framing と linking number の補正だけで，本論文では区別しない（KN の脚注と同じ扱い）．

\subsection{refined reformulated invariants}
KN は，$q$-brane と $\bar t$-brane の二つの配位について，幾何学的遷移の等式
\begin{align}
\sum_\lambda \overline{\rCS}_\lambda(K;a,q,t)\,g_\lambda(q,t)\,P_\lambda(x;q,t)
&=\exp\Bigl(\sum_{d=1}^\infty\sum_\mu \frac1d\,\frac{f^q_\mu(K;a^d,q^d,t^d)}{q^{\frac d2}-q^{-\frac d2}}\,s_\mu(x^d)\Bigr),\label{eq:GT1}\\
\sum_\lambda \overline{\rCS}_\lambda(K;a,q,t)\,P_{\lambda^T}(-x;t,q)
&=\exp\Bigl(\sum_{d=1}^\infty\sum_\mu \frac1d\,\frac{f^{\bar t}_\mu(K;a^d,q^d,t^d)}{t^{-\frac d2}-t^{\frac d2}}\,s_\mu(x^d)\Bigr)\label{eq:GT2}
\end{align}
を提案した．
ここから定まる refined reformulated invariants $f^q_\mu$, $f^{\bar t}_\mu$ は，行列 $M_{\mu\rho}(t)$ により
\begin{equation}\label{eq:f-to-hatf}
f^q_\mu(K;a,q,t)=\sum_\rho M_{\mu\rho}(t)\,\wh f_\rho(K;a,q,t),\qquad f^{\bar t}_\mu(K;a,q,t)=\sum_\rho M_{\mu\rho}(q^{-1})\,\wh f_\rho(K;a,q,t)
\end{equation}
と書け，二つの式が同じ $\wh f_\rho$ を与えることが KN の枠組みの整合性になる．
$f^q_\mu$ の尺度は KN の付録 B の式に合わせる．本文の \eqref{eq:GT1} をそのまま展開した値とは $(q/t)^{|\mu|/2}$ だけ違う（付録~\ref{app:typos}）．

KN の予想は，$a$ 次数を $\pm\frac12$ ずらした後で $\wh f_\rho$ が
\begin{equation}\label{eq:hat-f-2}
\wh f_\rho(K;a,q,t)=\sum_{g,\beta,J_r}(-1)^{2J_r}\,\wh N_{\rho,g,\beta,J_r}(K)\,(q^{\frac12}-q^{-\frac12})^{g}(t^{-\frac12}-t^{\frac12})^{g}\Bigl(\frac qt\Bigr)^{J_r-\frac\beta2}a^{\beta}
\end{equation}
の形をとり，$\wh N_{\rho,g,\beta,J_r}(K)$ が非負整数になる，というものである．
$a$ 次数のずらしは，$|\rho|=1$ で $a^{-1/2}$，$|\rho|=3$ で $a^{1/2}$ とする．
\eqref{eq:cov} の変数では，\eqref{eq:hat-f-2} は
\[
\wh f_\rho=\sum_{g,\beta,J_r}\wh N_{\rho,g,\beta,J_r}\,(\bq\bt-\bq^{-1}\bt^{-1})^{g}(\bq-\bq^{-1})^{g}\,\ba^{2\beta}\,\bt^{2J_r}
\]
となる（KN の $\wh{\mathbf N}_{\rho,g,\beta,F}$ で $F=2J_r$ としたもの）．
$G=(\bq\bt-\bq^{-1}\bt^{-1})(\bq-\bq^{-1})$ とおき，この展開を $G$ 展開と呼ぶ．
$G$ 展開が整数係数で閉じるのは，$\ba^{2\beta}$ の各行が
\begin{equation}\label{eq:sigma}
\sigma:\ (e_\ba,e_\bq,e_\bt)\mapsto(e_\ba,-e_\bq,e_\bt-e_\bq)
\end{equation}
で不変で，$\bq$ の指数が偶数のときである．
KN は，展開をする前の素朴な変数変換の係数
\[
\wh f_\rho(K;a=-\ba^2\bt,\,q=\bq^2\bt^2,\,t=\bq^2)=\sum_{i,j,k}\wh{\mathbf N}_{\rho;i,j,k}(K)\,\ba^{2i}\bq^{2j}\bt^{k}
\]
も非負整数になることを観察した．
本論文で refined LMOV 不変量と呼ぶのは，$G$ 展開の前の $\wh{\mathbf N}_{\rho;i,j,k}$ と，後の $\wh N_{\rho,g,\beta,J_r}$ の二つである．$f^q$, $f^{\bar t}$ の係数は付録で補足として扱う．

\begin{definition}[正値性]\label{def:positivity}
$\wh f_\rho$ が\emph{正}であるとは，全体の符号を一つ選ぶと，$\wh{\mathbf N}_{\rho;i,j,k}$ と $\wh N_{\rho,g,\beta,J_r}$ の全てが非負整数になることをいう．
\end{definition}

$\yng(1)$ では $\hf_{\yng(1)}=(a^{1/2}-a^{-1/2})\rCS_{\yng(1)}$ である．
したがって，$\rCS_{\yng(1)}$ を太字の変数で書いた $P$ に対し，$\hf_{\yng(1)}$ は（a 次数をずらした後）$(1+\ba^2\bt)\ba^{-2}\bt^{-1}P$ になる．
$P$ が HOMFLY-PT ホモロジーの Poincaré 多項式として解釈されるなら，$\yng(1)$ の正値性は自然に期待される．

\section{framed knot への拡張}\label{sec:framing}

\subsection{定義}
% [F1]
\begin{definition}\label{def:framed}
分割 $\lambda$ に対し
\[
\fr[\lambda]=\bigl(-(q/t)^{1/2}\bigr)^{|\lambda|}\,q^{n(\lambda')}\,t^{-n(\lambda)}
=(-1)^{|\lambda|}\prod_{s\in\lambda}q^{a(s)+\frac12}\,t^{-l(s)-\frac12}
\]
とおく．framing $x\in\mathbb{Z}$ の refined CS 不変量を $\rCS^{(x)}_\lambda=\fr[\lambda]^x\,\rCS_\lambda$ で定め，
これを \eqref{eq:GT1}，\eqref{eq:GT2} に代入して得る量を $f^{q,(x)}_\mu$，$f^{\bar t,(x)}_\mu$，$\hf^{(x)}_\rho$，$\hbN^{(x)}$，$\hN^{(x)}$ と書く．
$x=0$ が KN の場合である．
\end{definition}

$\fr[\lambda]$ は係数 $\pm1$ の単項式であり，plethystic 対数の中の Adams 作用 $\rCS_\lambda\mapsto\rCS_\lambda(a^k,q^k,t^k)$ では
$\fr[\lambda](q^k,t^k)^x$ が付く．符号 $(-1)^{|\lambda|}$ は $k$ によらない．
この扱いは Mariño--Vafa~\cite{Marino01} の例
$f^{(p)}_{\yng(2)}=q^pW_{\yng(2)}-\tfrac12\bigl(W_{\yng(1)}^2+(-1)^pW_{\yng(1)}(q^2,\lambda^2)\bigr)$ と同じである．
この定義は，著者が以前（2017--2018 年，未発表）に行った計算の定義と同じである．
その出力 9 件と全係数で一致することを確かめた（付録~\ref{app:data}）．

\subsection{framing の因子の由来}
% [F2]
\begin{proposition}\label{prop:MV}
$t=q$ で $\fr[\lambda]^x=(-1)^{|\lambda|x}q^{x\kappa_\lambda/2}$ である．
ここで $\kappa_\lambda=|\lambda|+\sum_i(\lambda_i^2-2i\lambda_i)=2\sum_{s\in\lambda}c(s)$ は Mariño--Vafa の $\kappa_R$ である．
したがって $\rCS^{(x)}$ の $t=q$ の極限は，Mariño--Vafa~\cite{Marino01} の framed knot の規則 $W^{(p)}_R=(-1)^{\ell p}q^{p\kappa_R/2}W_R$ に $x=p$ で一致する．
\end{proposition}
\begin{proof}
$t=q$ で $(q/t)^{1/2}=1$，$n(\lambda')-n(\lambda)=\sum_s c(s)=\kappa_\lambda/2$ である．
変数の向き（$q$ と $q^{-1}$）は，自明な結び目の主特殊化と Mariño--Vafa の量子次元を比べて確かめる．
$a=\lambda_{\mathrm{MV}}$，$q=q_{\mathrm{MV}}$ で両者は単項式倍を除いて一致し，$q\mapsto q^{-1}$ では一致しない（$\lambda\neq\lambda^T$ のとき）．
\end{proof}

\begin{remark}[refined topological vertex との関係]\label{rem:Taki}
Taki~\cite{Taki07} の refined topological vertex の framing の因子は
$f_\mu(t,q)=(-1)^{|\mu|}t^{\|\mu^T\|^2/2}q^{-\|\mu\|^2/2}$ である．
$\|\mu\|^2=2n(\mu')+|\mu|$ を使うと $f_\mu(t,q)=\fr[\mu]^{-1}$ となり，定義~\ref{def:framed} の因子に（向きを除いて）ちょうど一致する．
Taki は，Iqbal--Kozçaz--Vafa~\cite{Iqbal07} の因子 $(-1)^{|\nu|}t^{n(\nu)}q^{-n(\nu^T)}$ に $(t/q)^{|\nu|/2}$ を掛ける修正が，
K 理論版の Nekrasov 分配関数と一致するために必要であることを示した．
また Aganagic--Shakirov の refined $T$ 行列の固有値 $T_i=c_Tq^{(\lambda_{R_i}+\beta\rho)^2/2}$（$t=q^\beta$）について，
$T_\lambda/T_\emptyset$ は $SU(N)$ の部分を除いて $(a^{1/2}(q/t)^{1/2})^{|\lambda|}\,q^{n(\lambda')}t^{-n(\lambda)}$ に等しい．
箱あたりの単項式を除けば，framing の 1 単位は refined の $T$ 行列の作用である．
\end{remark}

% [F6]
\begin{proposition}\label{prop:bold}
太字の変数では
\[
\fr[\lambda]=\bt^{\,|\lambda|+2n(\lambda')}\,\bq^{\,2(n(\lambda')-n(\lambda))}
\]
であり，符号を含まない．例えば $\fr[\square]=\bt$，$\fr\yng(2)=\bq^2\bt^4$，$\fr\yng(1,1)=\bq^{-2}\bt^2$，$\fr\yng(2,1)=\bt^5$ である．
\end{proposition}
\begin{proof}
\eqref{eq:cov} から $-(q/t)^{1/2}=\bt$，$q=\bq^2\bt^2$，$t=\bq^2$ である．
\end{proof}

\begin{remark}
命題~\ref{prop:bold} により，framing の変化は（彩色付き）HOMFLY-PT ホモロジーの次数を符号なしでずらす操作として働く．
非 refined の極限 $\bt=-1$ では $\fr[\lambda]^x\to(-1)^{(|\lambda|+2n(\lambda'))x}=(-1)^{|\lambda|x}$ となる．
Mariño--Vafa の符号 $(-1)^{\ell p}$ は，refined の次数のずらし $\bt^{\ell p}$ の非 refined での影と読める．
この符号が KN の形に必要であることは命題~\ref{prop:sign} で示す．
\end{remark}

\subsection{$x$ についての構造}
% [F5]
\begin{lemma}\label{lem:adams}
$c$ を係数 1 の単項式とし，Adams 作用で $c\mapsto c^k$ となるものとする（$c=a^{1/2}$, $(q/t)^{1/2}$ など）．
$\rCS_\lambda$ を $c^{|\lambda|}\rCS_\lambda$ に置き換えると，$\hf_\rho$ は $c^{|\rho|}\hf_\rho$ になる．
\end{lemma}
\begin{proof}
$c^{|\lambda|}$ は生成関数の変数を $x\mapsto cx$ と取り替えることに当たり，$c$ が Adams 作用と両立するので plethystic 対数と可換である．
\end{proof}
太字の変数で $c=(q/t)^{1/2}=-\bt$ のとき，これは $\hbN$，$\hN$ の $2J_r$ を一様にずらし，全体の符号を変えるだけである．
したがって正値性に影響しない．

% [F7]
\begin{proposition}\label{prop:expsum}
$|\rho|=n$ のとき，$\bt^{-nx}\hf^{(x)}_\rho$ は $x$ の指数和である．根は
$n=2$ で $\{\bq^2\bt^2,\bq^{-2},1,-1\}$，$n=3$ で $\{\bq^6\bt^6,\bq^{-6},\bq^2\bt^2,\bq^{-2},\bt^2,1,-1\}$ の 7 個である．
$x$ の偶奇で分ければ，根は係数 1 の単項式になる．
\end{proposition}
\begin{proof}
$\hf_\rho$ は，$\rCS_\lambda$ とその Adams の写しの積を固定した有理関数の係数で一次結合したものである．
framing をつけると，各項には $\prod\fr[\lambda_i](q^{k_i},t^{k_i})^x$ が掛かる．
これらを太字の変数で書き出すと，$n=2$ では $\fr\yng(2)=\bq^2\bt^4$，$\fr\yng(1,1)=\bq^{-2}\bt^2$，$\fr[\square]^2=\bt^2$，$\fr[\square](q^2,t^2)=-\bt^2$ である．
$\bt^{-2x}$ を掛ければ主張の根が得られる．$n=3$ も同様である．
\end{proof}

指数和の形は，ツイスト結び目 $K_p$ の撚りの数 $p$ についての構造（\S\ref{sec:two}）と同じ型である．
framing は「もう一つの撚りの変数」として振る舞う．
$v=\bq^2\bt$ とおくと $\bq^2\bt^2=v\bt$，$\bq^{-2}=\bt/v$ で，根は $\sigma$ が入れ替える $v^{\pm1}$ の対をなす．

% [F8, F9]
\begin{proposition}\label{prop:x-indep}
任意の $x$ について次が成り立つ．
\begin{enumerate}
\item $\fr[\lambda](1/t,1/q)=\fr[\lambda^T](q,t)$．したがって，$\rCS_\lambda$ の転置対称性，$q$-brane と $\bar t$-brane の式 \eqref{eq:GT1}，\eqref{eq:GT2} が同じ $\hf$ を与えること，$G$ 展開が閉じる条件（$\sigma$ 不変性）は，$x$ によらない．
\item 2 箱について，$\hf^{(x)}_\rho$ が Laurent 多項式であることと $\hf^{(0)}_\rho$ が Laurent 多項式であることは同値である．
\end{enumerate}
\end{proposition}
\begin{proof}
(1) は定義から直接従う．
(2) $|\rho|=2$ の $\hf_\rho$ は，四つの入力 $\rCS_{\yng(2)}$，$\rCS_{\yng(1,1)}$，$\rCS_{\yng(1)}^2$，$\rCS_{\yng(1)}(a^2,q^2,t^2)$ の一次結合を，
共通の分母 $2A^2(Q^4-1)(T^4-1)(Q^2T^2-1)$（$(A,Q,T)=(a,q,t)^{1/2}$）で割ったものである．
分母の各既約因子の零点で分子に残る入力はちょうど二つで，その二つの framing の因子は零点の上で一致する（記号計算で確かめた）．
よって，零点での分子の消え方は $x$ によらない．
\end{proof}

% [F10]
\begin{proposition}\label{prop:sign}
定義~\ref{def:framed} の符号 $(-1)^{|\lambda|x}$ を外し，$T_\lambda^x$ だけを掛けたとする．
このとき $x$ が奇数なら，命題~\ref{prop:x-indep} (2) の分母の零点 $Q=i$，$T=i$（$q=-1$，$t=-1$）で残る二つの入力の因子が一致しない．
したがって，一般の入力で $\hf_{\yng(2)}$，$\hf_{\yng(1,1)}$ は Laurent 多項式でなくなる．
$x$ が偶数なら符号は効かない．また，箱あたりの $(q/t)^{1/2}$ を外しても Laurent 性は保たれる．
\end{proposition}
\begin{proof}
符号を外すと，Adams の写し $\rCS_{\yng(1)}(a^2,q^2,t^2)$ の項の因子だけが $(-1)^x$ 倍変わる．この項が残る零点で一致が壊れる．
\end{proof}

\begin{compresult}
三葉結び目，$x=1$ で符号を外すと，$\hf_{\yng(2)}$，$\hf_{\yng(1,1)}$ の分母に $(1+q)(1+t)$ が残る．
$x=2$ と，符号をつけた $x=1$ では Laurent 多項式である．
\end{compresult}

命題~\ref{prop:sign} は，Mariño--Vafa の「符号は整数性に不可欠である」という観察~\cite{Marino01} の refined 版である．
非 refined の場合，符号は平坦座標のずれ $\hat u\mapsto\hat u+i\pi p$ から説明された．refined の場合も，平坦座標のずれ $V\mapsto\bt V=-(q/t)^{1/2}V$ として読める．
Aganagic--Shakirov~\cite{Aganagic11} も，三葉結び目の前因子に現れる $\sqrt{q/t}$ が平坦座標のずれに関係するはずだと述べている．
ただし，このずれを幾何から導くことは本論文ではしていない．

\subsection{KN の命題の framed 版}
\begin{lemma}[$G$ 展開の補題]\label{lem:G}
$\ba^{2\beta}$ の各行が $\sigma$ で不変で $\bq$ の指数が偶数（すなわち $G$ 展開が閉じる）であり，展開の前の係数が非負なら，展開の後の係数 $\hN$ も非負整数である．
逆は成り立たない．
\end{lemma}
\begin{proof}
$y=\bq^2\bt$ とおくと，各行は $\bt^l(y^k+y^{-k})$ の非負結合である．
$p_k=y^k+y^{-k}$ を $w=y+y^{-1}$ の多項式とみなす．
$w_0=\bt+\bt^{-1}$ のまわりの Taylor 係数は $dp_k/dw=kU_{k-1}(w/2)$ と母関数
\[
\frac{1}{1-wx+x^2}=\sum_{g\ge0}\frac{G^g x^g}{\bigl((1-\bt x)(1-x/\bt)\bigr)^{g+1}}
\]
から非負である（$G=w-w_0$）．逆が成り立たない例は非代数的なケーブル $K5c$ にある（付録~\ref{app:nonalg}）．
\end{proof}

% [F12, F13]
\begin{proposition}\label{prop:one}
$P\in\mathbb{Z}_{\ge0}[\ba^{\pm1},\bq^{\pm1},\bt^{\pm1}]$ の $\bq$ の指数が偶数で，$P$ が $\sigma$ \eqref{eq:sigma} で不変とする．
$P$ を $\rCS_{\yng(1)}$ とする $\hf_{\yng(1)}$ は，全ての $x$ について，$G$ 展開が整数で閉じ，展開の前も後も非負である．
\end{proposition}
\begin{proof}
$x=0$ では，$\hf_{\yng(1)}=(1+\ba^2\bt)\ba^{-2}\bt^{-1}P$ も $\sigma$ で不変で非負なので，補題~\ref{lem:G} から従う．
一般の $x$ では $\fr[\square]^x=\bt^x$ なので，$\hf^{(x)}_{\yng(1)}=\bt^x\hf_{\yng(1)}$ である．これは $\hN$ の $2J_r$ を $x$ だけずらす．
\end{proof}

\begin{corollary}\label{cor:one}
$\rCS_{\yng(1)}$ が結び目の reduced HOMFLY-PT ホモロジーの Poincaré 多項式（DGR の正規化）で与えられるなら，任意の結び目・任意の $x$ で $\yng(1)$ は正である．
\end{corollary}
\begin{proof}
係数は次元なので非負である．$\sigma$ 不変性は Gorsky--Hogancamp--Mellit~\cite{Gorsky21} の対称性の定理である．
$\bq$ の指数が偶数であることは結び目のホモロジーの次数の偶奇から従う．
\end{proof}

% [F15]
\begin{proposition}\label{prop:two-necessary}
非負係数の入力 $\rCS_{\yng(1)}$，$\rCS_{\yng(2)}$ に対し，$\hf_{\yng(2)}$ が Laurent 多項式なら $\rCS_{\yng(2)}|_{q=1}=\rCS_{\yng(1)}|_{q=1}^2$ である．
特に，入力がホモロジーの Poincaré 多項式なら $\dim H^{\yng(2)}=(\dim H^{\yng(1)})^2$ が必要である．
$\yng(1,1)$ についても $t=1$ で同様である．この条件は $x$ によらない．
\end{proposition}
\begin{proof}
$\hf_{\yng(2)}$ の分子の $Q=1$ での値は $\rCS_{\yng(2)}-\rCS_{\yng(1)}^2$ に比例する．
非負係数の入力では，$q=1$（$\bt=-\bq^{-1}$）で各項の符号が $\bq$ の指数の偶奇だけで決まり，打ち消しが起きない．
したがって次元の等式が必要になる．$x$ によらないことは命題~\ref{prop:x-indep} による．
\end{proof}

\begin{remark}
$\mathbf{9_{42}}$ では $\dim H^{\yng(1)}=9$，$\dim H^{\yng(2)}=401$~\cite{Gorsky13} で $401\neq81$ である．
したがって，$t_c$ 次数の選び方によらず，KN の形に乗らない．
$\mathbf{9_{42}}$ が KN の形に乗らないこと自体は KN が述べている．命題~\ref{prop:two-necessary} はその理由を与える．
\end{remark}

\subsection{全ての framing での判定法}\label{sec:newton}
命題~\ref{prop:expsum} により，固定した結び目について全ての $x$ の正値性を有限個の判定に帰着できる．
$x$ を偶奇と向き（$x\ge x_0$ または $x\le x_0$）に分け，$x=x_0+2y$ について，根の並び $(r_0,r_1,r_2)$ のニュートン形
\[
\bt^{-nx}\hf^{(x)}_\rho=\sum_{k}c_k\,h_{y-k}(r_0,\dots,r_k)
\]
を作る（$h$ は完全斉次対称式）．係数 $c_k$ が全て同じ符号の非負係数の Laurent 多項式なら，$h$ の係数が非負なので，$y\ge0$ の全てで展開の前が非負である．
$\sigma$ 不変性は命題~\ref{prop:x-indep} から全ての $x$ に従い，補題~\ref{lem:G} で展開の後も非負になる．
起点 $x_0$ と 0 の間は直接判定する．

\subsection{正規化との関係}\label{sec:N2-framing}
% [F17]
\begin{proposition}\label{prop:N2-framing}
\S\ref{sec:normalization} の正規化 (N2$'$) の単項式 $M_\lambda=M_\square^{|\lambda|}T_\lambda^F$（$F=p_2q_2$）は
\[
M_\lambda=\bigl(M_\square(-(q/t)^{1/2})^{-F}\bigr)^{|\lambda|}\,\fr[\lambda]^F
\]
と書ける．すなわち，(N2$'$) のうち彩色によって変わる部分は，ちょうど $F$ 単位の framing である．
基準の framing を偶数 $y$ だけずらすことは，$x$ を $y$ ずらし，補題~\ref{lem:adams} の一様なずれを加えることに等しい．
奇数だけずらすと Mariño--Vafa の符号が落ち，命題~\ref{prop:sign} により Laurent 性が壊れる．
\end{proposition}
したがって，全ての $x\in\mathbb{Z}$ についての正値性の主張は，$x=0$ をどこに置くかに依存しない．
Aganagic--Shakirov~\cite{Aganagic11} が三葉結び目の値と DGR の superpolynomial の比を「$-6$ 単位の framing と $\sqrt{q/t}$」と書いていることは，$F=mn=6$ に合う（向きの規約を除く）．

\subsection{物理的な描像と予想}\label{sec:physics}
framed knot の refined large $N$ 双対は，次のように組み立てられる．
\begin{itemize}
\item Chern--Simons 側：framing の変化は refined CS 理論の $T$ 行列の作用である（注意~\ref{rem:Taki}）．
Aganagic--Shakirov の構成では，$U(1)_R$ 対称性は 3 次元多様体が半自由な $U(1)$ 作用をもつ（Seifert 多様体である）ことから存在する．
トーラス結び目の framing の変化は，同じ $S$, $T$ 行列の形式の中で扱われる．
\item 弦側：Mariño--Vafa の議論では，framing の整数は $L_K$ の平坦座標の取り方（無限遠での境界条件）に当たり，BPS 状態の数え方の構造は変わらない．
KN の Schwinger の計算と抽出式 \eqref{eq:GT1}，\eqref{eq:GT2} は，全ての $x$ でそのまま適用できる．
\end{itemize}
弦側で $U(1)_R$ の作用が自明なまま保たれることは，導出していない．
この点を仮定として，次を予想する．

% [F19]
\begin{tcolorbox}
\begin{conjecture}[framed refined LMOV 不変量の正値性]\label{conj:A}
トーラス結び目，および $\rCS_\lambda$ が彩色 HOMFLY-PT ホモロジーの Poincaré 多項式で与えられる結び目について，
全ての $x\in\mathbb{Z}$ と全ての $\rho$ で $\hf^{(x)}_\rho$ は正（定義~\ref{def:positivity}）である．
\end{conjecture}
\end{tcolorbox}

\section{入力：KN のデータを超える refined 不変量}\label{sec:inputs}

\subsection{トーラス結び目}
トーラス結び目 $T_{m,n}$ の $\rCS_\lambda$ は，Gorsky--Neguţ~\cite{Gorsky13b} による楕円 Hall 代数の生成元の作用で計算する．
楕円 Hall 代数と DAHA の対応は Schiffmann--Vasserot~\cite{Schiffmann09} により，ケーブルの場合の対応は Morton--Samuelson~\cite{Morton14} による．
% [I1]
規約を固定するため，KN の補助ファイルにあるトーラス結び目の値 27 件（$T_{2,3}$, $T_{2,5}$, $T_{2,7}$, $T_{2,9}$ の $|\lambda|\le3$ と $T_{3,4}$ の $\yng(1),\yng(2),\yng(1,1)$）と，
KN の付録 C の BPS の表 13 個を，単項式の正規化（\S\ref{sec:normalization}）の下で全係数で再現した．
これは既知の結果の再現であり，以下の計算の規約の確認として用いる．

\subsection{代数的な iterated torus 結び目}
companion $K$ の $(p,q)$-ケーブルを $C_{p,q}K$ と書く．$q$ は companion の Seifert framing で測る．
Cherednik--Danilenko~\cite{Cherednik14,Cherednik15} の DAHA による定義では，Newton 対は $r=(p_1,p_2)$，$s=(q_1,q_2-p_1q_1p_2)$ であり，
全ての $r_i,s_i>0$（$q_2>p_1q_1p_2$）のとき代数的である．
% [I2]
我々はこの定義を companion の段で線形に分解し，楕円 Hall 代数の演算子の合成として計算した．
すなわち $F_\lambda=P_\lambda[p_k\mapsto P_{k(p,-s)}]\cdot1=\sum_\mu c_{\lambda\mu}(s)P_\mu$ と展開し，
companion の値 $\rCS_K(\mu)$ との一次結合をとる．
素朴な Rosso--Jones の $t$ 変形が誤りであることは知られており~\cite{Danilenko14,Morozov18}，ここで使う形は定義から直接従うものである．
\begin{compresult}\label{cr:daha}
この計算法は，Cherednik--Danilenko の明示例 10 件（$K1$ の $\yng(1)$ と $\yng(2)$，$K2,K3,K4,K5a,K5b,K5c,K7,K8$ の $\yng(1)$）と，単項式倍を除いて全係数で一致する．
また，$t=q$ の極限はブレイドから計算した彩色 HOMFLY-PT 多項式と一致する．
\end{compresult}

本論文で扱うケーブルを表~\ref{tab:cables} にまとめる．
\begin{table}[h]
\centering
\caption{代数的な iterated torus 結び目．$b$ はブレイドの本数，$\delta$ は種数型の量（\S\ref{sec:normalization}）．}\label{tab:cables}
\begin{tabular}{llcc}
\toprule
名前 & 結び目 & $b$ & $\delta$ \\
\midrule
$K1$ & $T_{2,3}$ の $(2,13)$-ケーブル & 4 & 8 \\
$K3$ & $T_{2,3}$ の $(2,15)$-ケーブル & 4 & 9 \\
$K6a$ & $T_{2,5}$ の $(2,21)$-ケーブル & 4 & 14 \\
$K9$ & $T_{2,7}$ の $(2,31)$-ケーブル & 4 & 21 \\
$K8$ & $T_{3,4}$ の $(2,25)$-ケーブル & 6 & 18 \\
$K7$ & $T_{2,3}$ の $(3,19)$-ケーブル & 6 & 21 \\
$K10$ & $K1$ の $(2,53)$-ケーブル & 8 & 42 \\
\bottomrule
\end{tabular}
\end{table}

\subsection{単項式の正規化}\label{sec:normalization}
DAHA と楕円 Hall 代数は値を単項式倍を除いてしか決めないが，KN の BPS 抽出は単項式倍で結果が変わる．
そこで次の規則で正規化する．
\begin{itemize}
\item[(N1)] $t=q$ の極限が，framing 0 の reduced 彩色 HOMFLY-PT 多項式（DGR の規約）に等しい．
\item[(N2)] $\yng(1)$ の太字の変数での最低の $\bt$ 次数が 0 である．
\item[(N2$'$)] 彩色 $\lambda$ の単項式は $M_\lambda=M_\square^{|\lambda|}\,T_\lambda^{F}$ である．$F$ は一番外側のケーブル $C_{p,q}$ では $pq$，トーラス結び目 $T_{m,n}$ では $mn$ である．
\end{itemize}
トーラス結び目では，(N1)(N2) の下で KN の値 27 件が全係数で再現される．
(N2) が決めるのは，箱あたりの $(q/t)^{c}$ の自由度だけである．補題~\ref{lem:adams} により，これは正値性に影響しない．
彩色付きの単項式 $\pm a^\alpha q^\beta t^\gamma$ のうち，$t=q$ で決まるのは符号，$\alpha$，$\beta+\gamma$ である．
彩色付きの (N1) は検査として確かめた．
$t=q$ で決まらないのは，彩色ごとの $(q/t)^{c_\lambda}$ の自由度だけである．

% [I4]
\begin{hypothesis}\label{hyp:N2}
ケーブルについて，$t=q$ で決まらない彩色ごとの単項式 $(q/t)^{c_\lambda}$ は (N2$'$) の形をとる．
すなわち，$|\lambda|$ に比例する部分（正値性に無関係）と，refined の framing の固有値 $T_\lambda^F$ の分だけである．
\end{hypothesis}
命題~\ref{prop:N2-framing} により，仮説~\ref{hyp:N2} は「彩色によって変わる部分は Seifert framing の $F$ 単位の framing である」と言い換えられる．
根拠は次の三つである．
(i) 後述の refined 衛星公式（仮説~\ref{hyp:sat}）から従う．
(ii) 非代数的なケーブル $K5c$ では，$\yng(2)$ を $(q/t)^u$ 倍する 1 パラメータ族のうち，Laurent 性と転置対称性を満たすのは (N2$'$) の $u$ だけである．
(iii) 9 個のケーブルの全ての彩色で，(N2$'$) の形が計算した値と一致する．

% [I5]
\begin{compresult}\label{cr:M}
(N1)(N2) で決めた $\yng(1)$ の単項式は，保存した全 22 結び目（トーラス結び目 9 個，ケーブル 13 個）で
\[
M_\square=(-q/t)^{b-1}\bigl(a(q/t)^{1/2}\bigr)^{\delta}
\]
に一致する．ここで $b$ はブレイドの本数，$\delta$ は反復ごとに $\delta\mapsto p\delta+(p-1)(q-1)/2$ で決まる（$T_{m,n}$ では $\delta=(m-1)(n-1)/2$）．
\end{compresult}

% [I6]
\begin{hypothesis}[refined 衛星公式]\label{hyp:sat}
$\rCS_{C_{p,q}K}(\lambda)=\sum_\mu C_{\lambda\mu}\,\rCS_K(\mu)$ が成り立つ．
係数 $C$ は，トーラス結び目 $T_{p,q}$（自明な結び目のケーブル）の正規化された係数 $C=M_{T_{p,q}}^{|\lambda|}T_\lambda^{pq}c(q)$ である．
\end{hypothesis}
$t=q$ では仮説~\ref{hyp:sat} は古典的な衛星公式に帰着し，(N1) と両立する．

\begin{theorem}\label{thm:M}
仮説~\ref{hyp:sat} の下で，計算結果~\ref{cr:M} の閉じた式と (N2$'$) の形が成り立つ．
ただし，トーラス結び目の場合の $(q/t)$ の指数の導出には，Mellit~\cite{Mellit17} の定理（dyck paths）と，次の式を用いる．
\[
\mathcal{P}_{m,n}=-\bigl(a(qt)^{-1/2}\bigr)^{\delta}\frac{1}{1-t}\,\frac{(Q_{m,n}\cdot1)[a/q-1]}{a/q-1}
\]
ここで $Q_{m,n}$ は Bergeron--Garsia--Leven--Xin~\cite{Bergeron14} の演算子で，この式は計算で確かめたものである．
\end{theorem}
\begin{proof}[証明の概略]
ケーブルの部分では，framing 演算子 $\tau:P_\mu\mapsto T_\mu P_\mu$ が楕円 Hall 代数の傾きを $(p,n)\mapsto(p,n+kp)$ に送ることを使う．
これと仮説~\ref{hyp:sat} から，$M_{C_{p,q}K}=M_K^p\,M_{T_{p,q}}$ と $F=pq$ が従う．
トーラス結び目の部分は四つの事実からなる．
(I) Gorsky--Neguţ の正規化前の値の最低 $a$ 係数は，有理 $q,t$-Catalan 数で与えられる．これは Gorsky--Neguţ の予想であり，Mellit~\cite{Mellit16} の定理の系として成り立つ．
(II) 最低 $a$ 次数は，Franks--Williams の定理による．
(III) $t=q$ での符号と $q$ の冪は，$q\leftrightarrow q^{-1}$ 対称性から決まる．
(IV) $(q/t)$ の指数は，Mellit の Dyck 路の表示から決まる．
詳細は付録~\ref{app:normalization} に述べる．
\end{proof}

\begin{remark}[Mellit の系の正規化]
Mellit~\cite[Cor.~daha]{Mellit17} は $\mathcal{P}_{m,n}=(-1)^n(a(qt)^{-1/2})^\delta\frac{1}{1-t}\frac{P_{m,n}(1)[a-1]}{a-1}$ と書き，$P_{m,n}$ を「対応する楕円 Hall 代数の生成元」とする．
同論文の Thm.~daha から，この式は $P_{m,n}(1)=-d_-B_{c_{m,n}}d_+(1)$（Carlsson--Mellit の $d_\pm$ の表現の元）と読めば成り立つ．
一方，標準の正規化（Bergeron--Garsia--Leven--Xin の $Q_{m,n}$，Gorsky--Neguţ の $\widetilde P$）の元と読むと，代入 $[a-1]$ では成り立たない．
候補の総当たりで成り立つのは定理~\ref{thm:M} の形だけである．
定理~\ref{thm:M} の (IV) が決めるのは $M_\square$ の $(q/t)$ の指数だけなので，この点は正値性の結果に影響しない（補題~\ref{lem:adams}）．
\end{remark}

\begin{remark}[Cherednik--Danilenko の tilde 正規化との関係]\label{rem:tilde}
Cherednik--Danilenko~\cite{Cherednik14} は DAHA superpolynomial を，最低の $a$ の冪と $q_{st}^\bullet t_{st}^\bullet$ で割って，$a_{st}=0$ の部分の定数項が 1 になるように正規化する（tilde 正規化）．
標準の変数 $(q_{st},t_{st},a_{st})$ は太字の変数と同じである．
係数が非負の保存値 19 個（トーラス結び目 9 個，ケーブル 10 個）で，本論文の $\yng(1)$ の太字の $P$ は
\[
P=\ba^{2\delta}\bq^{-2\delta}\times(\text{tilde 正規化})
\]
である（計算）．Cherednik--Danilenko の明示例との比も，正のもの（$K1,K3,K5b,K7,K8$）で全て $\ba^{2\delta}\bq^{-2\delta}$ である．
これらの $P$ は全て $\sigma$ で不変で，$\bq$ の指数は偶数である．
\end{remark}

\subsection{ホモロジーを入力とする結び目}
11 交点以下の全素結び目 801 個について，reduced HOMFLY-PT ホモロジーの Poincaré 多項式（Nakagane--Sano~\cite{Nakagane21} のデータ）を $\yng(1)$ の入力とする．
DGR の規約への変換は，$\sigma$ 不変性，$\bt=-1$ での HOMFLY-PT 多項式，トーラス結び目の $\rCS_{\yng(1)}$ との一致で確かめた．

\subsection{2 箱の予想式}
双曲結び目の refined な $\yng(2)$, $\yng(1,1)$ は定義どおりには得られない．そこで文献の予想式を入力とする．
\begin{itemize}
\item ダブルツイスト結び目 $K_{m,n}$ とツイスト結び目 $K_p=K_{p,1}$：Kameyama--Nawata--Tao--Zhang（KNTZ）~\cite{Kameyama19} の，長方形の彩色 HOMFLY-PT ホモロジーの Poincaré 多項式の式（refined 版，予想）を用いる．
KNTZ の $K_{m,n}$ は，GNSSS~\cite{Gukov15} の $(2m-1,1,2n-1)$-pretzel 結び目 $K_{s,t}$（$s=m$，$t=n$）と同じ結び目である．
KNTZ の式は GNSSS の四重次数の予想式と一致する（$\yng(1)$，$\yng(2)$，$\yng(1,1)$ について．下の計算結果~\ref{cr:kntz}）．
$t=q$ では KNTZ の式は Fang--Zhou~\cite{Fang26} により証明されている．
KN の $t_c$ 次数への変換は，Gorsky--Gukov--Stošić~\cite{Gorsky13} の次数の関係 $\mathbf{Q}\mapsto\bq^R$，$\mathbf{t}_r\mapsto\bq^{-1}$，$\mathbf{t}_c\mapsto\bt\bq$（$R$ は長方形の行数）による．
変換は，三葉結び目 $K_{1,1}=T_{2,3}$ の $\yng(1)$，$\yng(2)$，$\yng(1,1)$，$\yng(2,2)$ が楕円 Hall 代数の値と全係数一致することで較正した（Nawata ら~\cite{Nawata12} も参照）．
\item thick な結び目 $\mathbf{10_{139}}$，$\mathbf{10_{152}}$：Stošić~\cite{Stosic24} の quiver 型の式と，KRSS~\cite{Kucharski17b} の予想式で $t_c$ 次数を作る．
この方法は，$\mathbf{8_{19}}=T_{3,4}$ で楕円 Hall 代数の値と $\yng(2)$ から $\yng(1,1,1)$ まで全項一致する．
\end{itemize}
\begin{compresult}\label{cr:kntz}
KNTZ の refined の式と GNSSS の四重次数の式は，$K_{m,n}$ の $\yng(1)$，$\yng(2)$ と，ミラー対称性 $\mathbf{t}_r\leftrightarrow\mathbf{t}_c$ を介した $\yng(1,1)$ で，
符号の四つの象限それぞれの $|m|,|n|\in\{1,2,3\}$ で全係数一致する（144 件）．
両方の式は $m,n$ のそれぞれについて（象限ごとに）根が 3 個以下の指数和なので，全ての $(m,n)$ で一致する．
また KNTZ の式の実装は，KNTZ の補助ファイルの値と全係数一致する．
\end{compresult}
これらを入力とする主張は「予想式を入力とした場合の正値性」であり，定義どおりの入力による 3 箱の結果（\S\ref{sec:three}）とは確からしさが異なる．

\subsection{厳密な計算}
$\rCS_\lambda$ と $\hf_\rho$ は，有理数の点ではなく，31 ビットの素数を法とする評価と補間で求める．
係数が大きいときは 2〜3 個の素数と中国剰余定理を用いる．
得られた多項式は，補間に使っていない素数の点で検証する．
最大の例 $K10$ の 3 箱は，$G$ 展開の前で約 183 万項，係数の最大は約 $1.1\times10^{19}$ である．
計算時間とメモリは付録~\ref{app:computation} にまとめる．

\section{検証 I：$[1]$}\label{sec:one}

系~\ref{cor:one} により，ホモロジーを入力とする限り，任意の結び目・任意の framing で $[1]$ は正である．
この節では，それを結び目表で確かめ，$2J_r$ の偶奇の振る舞いを調べる．

% [R1]
\begin{compresult}\label{cr:table}
11 交点以下の全素結び目 801 個（thin で交代的なもの 562 個，thin で非交代的なもの 143 個，thick なもの 96 個）で，
$\hf_{[1]}$ の $G$ 展開は整数で閉じ，展開の前も後も非負である．
\end{compresult}

ケーブルの $[1]$ については，注意~\ref{rem:tilde} のとおり，計算した値が非負で $\sigma$ 不変である．したがって命題~\ref{prop:one} を直接適用できる．
一般の代数的な iterated torus 結び目については，Cherednik--Danilenko の予想（\S\ref{sec:conjB}）の下で同じ結論が従う．

\subsection{$2J_r$ の偶奇と parity knot}
KN は，トーラス結び目では $(\rho,g,\beta)$ を固定すると $2J_r$ の偶奇がそろうことを観察し，これを $S_R$ の自明な作用の証拠とした．
$[1]$ ではこの条件を入力の言葉で書ける．

% [F14]
\begin{proposition}\label{prop:parity}
命題~\ref{prop:one} の仮定の下で，$\hN_{[1],g,\beta,J_r}$ の $2J_r$ の偶奇が $(g,\beta)$ ごとにそろうことは，次と同値である．
$P$ の $\ba$ の次数が連続する範囲ごとに，$e_\bt+e_\ba/2$ の偶奇が一定である．
$\ba$ の次数の台に隙間が無いとき，これは Chandler--Gorsky~\cite{Chandler22} の parity knot の条件（HHH 全体で $t_{\mathrm{DGR}}-a/2$ の偶奇が一定）と同じである．
\end{proposition}
\begin{proof}
$\hf_{[1]}$ の $\ba^{2\beta}$ の行には $P$ の隣り合う二つの行だけが入る．
$p_k$ の $G^g$ の係数に現れる $\bt$ の指数は $k-g$ と同じ偶奇をもつので，$2J_r\equiv e_\bt-g$ である．
$\ba$ の次数は偶数なので，$\pm e_\ba/2$ の偶奇は同じである．
\end{proof}

\begin{compresult}
801 個の全てで $\ba$ の次数の台に隙間は無く，$2J_r$ の偶奇がそろう結び目は 9 個である．
$\mathbf{3_1},\mathbf{5_1},\mathbf{7_1},\mathbf{9_1},\mathbf{11a_{367}},\mathbf{8_{19}},\mathbf{10_{124}}$（トーラス結び目）と，$\mathbf{10_{125}}$（thick），$\mathbf{10_{140}}$（thin）である．
これらは Chandler--Gorsky の定義による parity knot と一致する．
\end{compresult}

したがって，$2J_r$ の偶奇の一致は $S_R$ の作用の自明さそのものではなく，ホモロジーの parity の性質を反映している．
非トーラス結び目にも偶奇がそろう例がある．
parity knot の条件そのものは Chandler--Gorsky による．BPS 不変量の偶奇との同値性と，11 交点以下の一覧は本論文で与える．

\section{検証 II：2 箱}\label{sec:two}

この節の双曲結び目の入力は予想式（\S\ref{sec:inputs}）であり，主張は全て「その予想式を入力とした場合」のものである．
トーラス結び目と代数的なケーブルの入力は定義どおりの refined 不変量である．

\subsection{ツイスト結び目}
$K_p=K_{p,1}$ を KNTZ~\cite{Kameyama19} の規約のツイスト結び目とする（$K_1=\mathbf{3_1}$，$K_2=\mathbf{5_2}$，$K_{-1}=\mathbf{4_1}$，$K_{-2}=\mathbf{6_1}$．GNSSS~\cite{Gukov15} の規約と同じ）．
入力は KNTZ の refined の式である（\S\ref{sec:inputs}）．
$X=\ba^2\bt^2$，$Y=\ba^4\bq^4\bt^8$，$\bar Y=\ba^4\bq^{-4}\bt^4$ とおく．

\begin{lemma}\label{lem:twist-expsum}
$\hf_\rho(K_p)$（$|\rho|=2$）は $p$ について根 $\{1,X,X^2,Y,\bar Y\}$ の指数和である（$p<0$ では根を逆数にし，全体に $\bt^2$ が掛かる）．
\end{lemma}
\begin{proof}
KNTZ の式では，$p$ は捻りの因子 $\mathbf{Tw}^{(p)}_\lambda$（$\lambda\subseteq\yng(2)$ または $\lambda\subseteq\yng(1,1)$）に固有値の $p$ 乗としてだけ入る．
太字の変数では，これは $\bigl((\ba\bt)^{2j}(\bt\bq)^{2j(j-1)}\bigr)^p$（$j=0,1,2$）の形である（GNSSS の式でも同じ形）．
入力の根は $P_{\yng(1)}$ で $\{1,X\}$，$P_{\yng(2)}$ で $\{1,X,Y\}$，$P_{\yng(1,1)}$ で $\{1,X,\bar Y\}$ である．
$\hf$ は入力の積と Adams 作用の一次結合を，命題~\ref{prop:two-necessary} の分母で割ったものである．
\end{proof}

% [R2]
\begin{theorem}\label{thm:twist}
KNTZ の refined の式を入力とすると，全ての $p\neq0$ で $\hf_{\yng(2)}(K_p)$，$\hf_{\yng(1,1)}(K_p)$ は正である．
\end{theorem}
\begin{proof}[証明の概略]
並び $(1,X,X^2,Y,\bar Y)$ のニュートン形 $s_{1+n}=\sum_{k=0}^{4}c_k\,h_{n-k}(m_0,\dots,m_k)$ を作る．
$c_k$ は $s_1,\dots,s_5$（$|p|=1,\dots,5$ の値）の整数係数の一次結合なので，Laurent 多項式である．
両方の $\rho$，両方の $p$ の符号で，$c_0,\dots,c_4$ の全係数が同じ符号であることを有限個の厳密な計算で確かめる．
$h$ の係数は非負なので，全ての $p$ で展開の前が一様な符号になる．
$\sigma$ は根の集合を保つ（$Y\leftrightarrow\bar Y$）ので，$\sigma$ 不変性は $|p|\le5$ から全ての $p$ に従い，補題~\ref{lem:G} で展開の後も非負になる．
$c_2$ の正値性は構造からは導けていない．
しかし，根の部分集合ごとの完全斉次対称式の非負結合（線形計画法による証明書，有理数で厳密に検証）でも同じ結論が得られる．
この証明は $c_2$ を使わない．
\end{proof}

\begin{compresult}
$|p|\le30$ の 120 件で，spin 表示と genus 表示の両方で負の項は 0 である．
閉じた形は，作るのに使っていない $|p|=6,\dots,12,15,20,30$ で直接計算と全係数一致する．
\end{compresult}

\begin{remark}[単峰性]
$v=\bq^2\bt$ とおくと，$c_3$ と $c_4$ は構造から正であり，$\hf(K_p)$ は全ての $p$ で $v$ についての $\mathfrak{sl}_2$ 指標の正係数の和（単峰）になる（$c_0,\dots,c_2$ は有限個の判定）．
$\yng(1)$ の単峰性は Gorsky--Hogancamp--Mellit~\cite{Gorsky21} の hard Lefschetz の性質から全ての結び目で従う．
2 箱の単峰性の構造的な理由は未解明である．
\end{remark}

\subsection{ダブルツイスト結び目}
% [R4]
\begin{theorem}\label{thm:pretzel}
KNTZ の refined の式を入力とすると，ダブルツイスト結び目 $K_{m,n}$ の 2 箱は次の範囲で正である．
\begin{enumerate}
\item $m,n$ が同符号の全ての $(m,n)$．
\item $m<0<n$ で $\min(|m|,n)\le20$ の全ての $(m,n)$．
\end{enumerate}
\end{theorem}
\begin{proof}[証明の概略]
$\hf(K_{m,n})=\bt^{2(1-\mathrm{sgn})}F(m,n)$ と書くと，$F$ は $\mathbb{Z}^2$ 全体で一つの指数和である．
同符号の領域では，2 変数のニュートン形の係数が全て同じ符号になる．
異符号の領域では 2 変数のニュートン形は一様にならない．そこで直線ごとの 1 変数のニュートン形で $\min(|m|,n)\le20$ を示す．
\end{proof}

\begin{compresult}
$|m|,|n|\le20$ の 820 点と，$21\le|m|\le30$・$21\le n\le40$ の 200 点で，両方の $\rho$ についてspin 表示と genus 表示の両方で負の項は 0 である．
\end{compresult}

異符号で $|m|,n>20$ の領域は，くさび形への分割と鏡像型の恒等式で調べる．
くさびの方向のニュートン係数 $d_2,d_3,d_4$ の非正値性は，線形計画法の証明書で示した．
残りは，1 変数の三つの数列の非正値性に帰着する．三つとは，対角線 $F(-u,u)$ と，両方の $\rho$ の $d_1(u)$ である．
これらは $u\le25$ で成り立つが，一般には未証明である．

\subsection{thick な非トーラス結び目}
% [R5]
\begin{compresult}\label{cr:thick}
Stošić の quiver と KRSS の予想式から作った $\mathbf{10_{139}}$，$\mathbf{10_{152}}$ の $\yng(2)$，$\yng(1,1)$ の入力で，2 箱は正である．
次元は $441=21^2$，$841=29^2$ で，命題~\ref{prop:two-necessary} の必要条件を満たす．
$t=q$ の極限は，$(2,s)$-ケーブルから逆に解いた $\yng(2)$，$\yng(1,1)$ の彩色 HOMFLY-PT 多項式と一致する．
\end{compresult}
これらは，命題~\ref{prop:two-necessary} の必要条件を満たす thick な非トーラス結び目で 2 箱の正値性が成り立つ例である．
knots--quivers 対応の DT 不変量の正値性からは従わない．
対称表現だけの quiver 級数からは，$\yng(1,1)$ の入力を要する $\hf_{\yng(2)}$，$\hf_{\yng(1,1)}$ は出てこない．

\begin{remark}[正値性は自動的ではない]
$\mathbf{5_2}$ で $t_c$ 次数の割り当てを別の 12 通りに変えると，いずれも負の項が出る．
正値性は入力の細かい次数に依存する，非自明な性質である．
\end{remark}

\subsection{全ての framing}
% [R7]
\begin{theorem}\label{thm:two-all-x}
自明な結び目，$T_{2,3}$，$T_{2,5}$，$T_{2,7}$，$T_{3,4}$，$\mathbf{4_1}$，$K1$，$K3$ の 2 箱は，全ての $x\in\mathbb{Z}$ で正である．
$\mathbf{4_1}$ の入力は KN のデータである．
\end{theorem}
\begin{proof}
\S\ref{sec:newton} の方法による．各結び目・各 $\rho$・$x$ の偶奇と向きの四つの組で，ニュートン形の係数が一様な符号になる起点 $x_0$ がある．
起点と 0 の間の $x$ は直接判定した．
\end{proof}
負の向きの起点は，結び目が複雑になるほど 0 から遠い（$T_{2,2k+1}$ で $-(2k+2)$，$K1$ で $-20$，$K3$ で $-22$）．この意味は未解明である．

% [R8]
\begin{theorem}\label{thm:twist-x}
KNTZ の refined の式を入力とすると，ツイスト結び目の 2 箱は，$p$ と $x$ が同符号（$p>0,x\ge0$ または $p<0,x\le0$）の全ての $(p,x)$ で正である．
\end{theorem}
\begin{proof}[証明の概略]
$\bt^{-2x}\hf^{(x)}_\rho(K_p)$ は $p$ と $x$ の両方について指数和である（補題~\ref{lem:twist-expsum}，命題~\ref{prop:expsum}）．
テンソル積のニュートン形の係数を $|p|=1,\dots,5$ と 3 点の $x$ から作り，同符号の向きで一様な符号であることを確かめる．
\end{proof}

\begin{compresult}
逆符号の領域では，$|p|\le6$，$|x|\le14$ の 696 件で負の項は 0 である．この領域の証明書は得られていない．
\end{compresult}

\begin{compresult}[framing を振った直接の判定]\label{cr:framing-batch}
次の 1821 件（$(x,\rho)$ の組）で，spin 表示と genus 表示の両方で負の項は 0 である（9 件は $\hf=0$）．
3 箱までは，自明な結び目と $T_{2,3}$（$|x|\le6$），$\mathbf{4_1}$ と $T_{2,5}$（$|x|\le4$），$T_{2,7}$ と $T_{3,4}$（$|x|\le3$），
$T_{3,5},T_{2,11},T_{3,7},T_{4,5}$ と $K6a,K7,K8,K9$（$|x|\le10$），$K3$（$|x|\le2$），$K1$（$|x|\le2$ と $x=\pm6$）である．
$K10$ は 2 箱まで $|x|\le10$，3 箱は $|x|\le2$ である．
ツイスト結び目 $K_{\pm1},\dots,K_{\pm6}$（$|x|\le4$）は 2 箱までである．
最大の例は $K10$ の 3 箱，$x=2$ で，spin 表示で約 189 万項である．
\end{compresult}

\begin{table}[h]
\centering
\caption{framing を振った判定．$|\rho|=1,2,3$ の列は調べた framing $x$ の範囲，件数は判定した $(x,\rho)$ の組の数（「各」は 1 結び目あたり．$K10$ の 3 箱の $x=0$ は表~\ref{tab:three} の判定で，件数に含めない），負（spin/genus）は spin 表示と genus 表示の負の項の数の合計（全体の符号は良い方）．自明な結び目の 9 件は $\wh f=0$．}\label{tab:framing}
\resizebox{\textwidth}{!}{% scripts/paper_framing_table.py が logs/framing_batch.log と data/framed/ から生成。
% 判定件数の合計 1821（重なりは 1 件と数える）。
\begin{tabular}{llcccc}
\hline
結び目 & 箱の数: $x$ の範囲 & 件数 & 前の負 & 後の負 & $\wh f=0$ \\
\hline
自明な結び目 & 1: $-6..6$; 2: $-6..6$; 3: $-6..6$ & 78 & 0 & 0 & 9 \\
$T_{2,3}$ & 1: $-6..6$; 2: $-6..6$; 3: $-6..6$ & 78 & 0 & 0 & 0 \\
$T_{2,5}$ & 1: $-4..4$; 2: $-4..4$; 3: $-4..4$ & 54 & 0 & 0 & 0 \\
$T_{2,7}$ & 1: $-3..3$; 2: $-3..3$; 3: $-3..3$ & 42 & 0 & 0 & 0 \\
$T_{3,4}$ & 1: $-3..3$; 2: $-3..3$; 3: $-3..3$ & 42 & 0 & 0 & 0 \\
$T_{3,5}$ & 1: $-10..10$; 2: $-10..10$; 3: $-10..10$ & 126 & 0 & 0 & 0 \\
$T_{2,11}$ & 1: $-10..10$; 2: $-10..10$; 3: $-10..10$ & 126 & 0 & 0 & 0 \\
$T_{3,7}$ & 1: $-10..10$; 2: $-10..10$; 3: $-10..10$ & 126 & 0 & 0 & 0 \\
$T_{4,5}$ & 1: $-10..10$; 2: $-10..10$; 3: $-10..10$ & 126 & 0 & 0 & 0 \\
$\mathbf{4_1}$ & 1: $-4..4$; 2: $-4..4$; 3: $-4..4$ & 54 & 0 & 0 & 0 \\
$K1$ & 1: $-2..2$; 2: $-2..2$; 3: $-6,-2,-1,0,1,2,6$ & 36 & 0 & 0 & 0 \\
$K3$ & 1: $-2..2$; 2: $-2..2$; 3: $-2..2$ & 30 & 0 & 0 & 0 \\
$K6a$ & 1: $-10..10$; 2: $-10..10$; 3: $-10..10$ & 126 & 0 & 0 & 0 \\
$K7$ & 1: $-10..10$; 2: $-10..10$; 3: $-10..10$ & 126 & 0 & 0 & 0 \\
$K8$ & 1: $-10..10$; 2: $-10..10$; 3: $-10..10$ & 126 & 0 & 0 & 0 \\
$K9$ & 1: $-10..10$; 2: $-10..10$; 3: $-10..10$ & 126 & 0 & 0 & 0 \\
$K10$ & 1: $-10..10$; 2: $-10..10$; 3: $-2,-1,1,2$ & 75 & 0 & 0 & 0 \\
$K_{1}$ & 1: $-4..4$; 2: $-4..4$ & 27 & 0 & 0 & 0 \\
$K_{-1}$ & 1: $-4..4$; 2: $-4..4$ & 27 & 0 & 0 & 0 \\
$K_{2}$ & 1: $-4..4$; 2: $-4..4$ & 27 & 0 & 0 & 0 \\
$K_{-2}$ & 1: $-4..4$; 2: $-4..4$ & 27 & 0 & 0 & 0 \\
$K_{3}$ & 1: $-4..4$; 2: $-4..4$ & 27 & 0 & 0 & 0 \\
$K_{-3}$ & 1: $-4..4$; 2: $-4..4$ & 27 & 0 & 0 & 0 \\
$K_{4}$ & 1: $-4..4$; 2: $-4..4$ & 27 & 0 & 0 & 0 \\
$K_{-4}$ & 1: $-4..4$; 2: $-4..4$ & 27 & 0 & 0 & 0 \\
$K_{5}$ & 1: $-4..4$; 2: $-4..4$ & 27 & 0 & 0 & 0 \\
$K_{-5}$ & 1: $-4..4$; 2: $-4..4$ & 27 & 0 & 0 & 0 \\
$K_{6}$ & 1: $-4..4$; 2: $-4..4$ & 27 & 0 & 0 & 0 \\
$K_{-6}$ & 1: $-4..4$; 2: $-4..4$ & 27 & 0 & 0 & 0 \\
\hline
\end{tabular}
}
\end{table}

\section{検証 III：3 箱}\label{sec:three}

3 箱には非長方形の彩色 $\yng(2,1)$ が要る．
定義どおりの $\rCS_{\yng(2,1)}$ が得られるのは，トーラス結び目と iterated torus 結び目（DAHA，楕円 Hall 代数）である．
この節の主な結果は，これらの結び目についてのものである．

\subsection{トーラス結び目と代数的な iterated torus 結び目}
% [R10, R11]
\begin{compresult}\label{cr:three}
表~\ref{tab:three} の 15 個の結び目で，$|\rho|\le3$ の全 6 つの $\rho$ について $\hf_\rho$ は正（framing 0）である．
内訳はトーラス結び目 7 個，代数的な iterated torus 結び目 7 個，$\mathbf{4_1}$ である．
同じ全体の符号で，$G$ 展開の前の $\hbN$ と後の $\hN$ がともに非負整数になる．
\end{compresult}

\begin{table}[h]
\centering
\caption[3 箱までの判定（framing 0）]{3 箱までの判定（framing 0）．$\checkmark(n)$ は $G$ 展開の前後とも負の項が 0 で，$n$ は展開の前の項数．
$T_{2,3}$，$T_{2,5}$ と $\mathbf{4_1}$ は KN の表の再現である．$\mathbf{4_1}$ の $\yng(2,1)$ は KN の候補を用いた．}\label{tab:three}
\resizebox{\textwidth}{!}{% scripts/paper_summary.py が生成。\checkmark(n) は spin 表示と genus 表示の両方で負の項 0（n は spin 表示の項数）。それ以外は spin 表示の負/項数; genus 表示の負/項数（前後それぞれで全体の符号は良い方）。framing 0。
\begin{tabular}{llcccccc}
\hline
結び目 & 分類 & $\yng(1)$ & $\yng(2)$ & $\yng(1,1)$ & $\yng(3)$ & $\yng(2,1)$ & $\yng(1,1,1)$ \\
\hline
$T_{2,3}$ & torus & \checkmark\,(6) & \checkmark\,(26) & \checkmark\,(41) & \checkmark\,(128) & \checkmark\,(236) & \checkmark\,(282) \\
$T_{2,5}$ & torus & \checkmark\,(10) & \checkmark\,(112) & \checkmark\,(144) & \checkmark\,(753) & \checkmark\,(986) & \checkmark\,(1083) \\
$T_{3,4}$ & torus & \checkmark\,(21) & \checkmark\,(429) & \checkmark\,(504) & \checkmark\,(3249) & \checkmark\,(3802) & \checkmark\,(4053) \\
$T_{3,5}$ & torus & \checkmark\,(30) & \checkmark\,(766) & \checkmark\,(867) & \checkmark\,(5831) & \checkmark\,(6559) & \checkmark\,(6890) \\
$T_{2,11}$ & torus & \checkmark\,(22) & \checkmark\,(733) & \checkmark\,(813) & \checkmark\,(5667) & \checkmark\,(6260) & \checkmark\,(6519) \\
$T_{3,7}$ & torus & \checkmark\,(53) & \checkmark\,(1797) & \checkmark\,(1953) & \checkmark\,(13600) & \checkmark\,(14696) & \checkmark\,(15199) \\
$T_{4,5}$ & torus & \checkmark\,(68) & \checkmark\,(2356) & \checkmark\,(2558) & \checkmark\,(18729) & \checkmark\,(20195) & \checkmark\,(20883) \\
K1 & alg & \checkmark\,(104) & \checkmark\,(4141) & \checkmark\,(4414) & \checkmark\,(32732) & \checkmark\,(34655) & \checkmark\,(35558) \\
K3 & alg & \checkmark\,(122) & \checkmark\,(5178) & \checkmark\,(5483) & \checkmark\,(41138) & \checkmark\,(43289) & \checkmark\,(44300) \\
K6a & alg & \checkmark\,(275) & \checkmark\,(12784) & \checkmark\,(13267) & \checkmark\,(100115) & \checkmark\,(103458) & \checkmark\,(105035) \\
K9 & alg & \checkmark\,(564) & \checkmark\,(28584) & \checkmark\,(29309) & \checkmark\,(224004) & \checkmark\,(228995) & \checkmark\,(231354) \\
K8 & alg & \checkmark\,(617) & \checkmark\,(31499) & \checkmark\,(32427) & \checkmark\,(251627) & \checkmark\,(258022) & \checkmark\,(261098) \\
K7 & alg & \checkmark\,(834) & \checkmark\,(43533) & \checkmark\,(44631) & \checkmark\,(343739) & \checkmark\,(351218) & \checkmark\,(354818) \\
K10 & alg & \checkmark\,(4119) & \checkmark\,(227727) & \checkmark\,(230642) & \checkmark\,(1805429) & \checkmark\,(1825082) & \checkmark\,(1834640) \\
$4_1$ & kn & \checkmark\,(10) & \checkmark\,(42) & \checkmark\,(42) & \checkmark\,(223) & \checkmark\,(264) & \checkmark\,(223) \\
\hline
\end{tabular}
}
\end{table}

$T_{2,3}$，$T_{2,5}$ と $\mathbf{4_1}$ の結果は，KN の表の再現である．
KN の表に無いトーラス結び目 5 個（$T_{3,4}$ の 3 箱，$T_{3,5}$，$T_{2,11}$，$T_{3,7}$，$T_{4,5}$）は新しい表である．
ただし KN はトーラス結び目全般について予想しているので，これは予想の確認に当たる．
代数的な iterated torus 結び目 7 個の 3 箱は，本論文の調べた範囲では文献に無い．
3 本ケーブル $K7$ と 2 段の反復 $K10$ では，$|\lambda|=3$ の refined 不変量そのものも文献に見当たらない．
Cherednik--Danilenko~\cite{Cherednik14} は，$K7$ の $\yng(1)$ と $K10$ の DAHA-Betti 多項式を与えている．

% [R12]
\begin{remark}
入力の $\rCS_{\yng(2,1)}$ 自体は負の係数をもつが，$\hf_{\yng(2,1)}$ は正になる．
したがって，3 箱の正値性は入力の係数の正値性から直接は従わず，KN の抽出（plethystic な対数と $M$ 行列）を経て初めて現れる．
\end{remark}

\subsection{framing を変えた 3 箱}
% [R13]
\begin{compresult}\label{cr:three-x}
表~\ref{tab:three} の全ての結び目について，3 箱を $x\neq0$ でも判定した．
トーラス結び目 $T_{3,5},T_{2,11},T_{3,7},T_{4,5}$ と代数的なケーブル $K6a,K7,K8,K9$ は $|x|\le10$，$K10$ は $|x|\le2$，$K1$ と $K3$ は $|x|\le2$（$K1$ は $x=\pm6$ も），
自明な結び目と $T_{2,3}$ は $|x|\le6$，$T_{2,5}$ と $\mathbf{4_1}$ は $|x|\le4$，$T_{2,7}$ と $T_{3,4}$ は $|x|\le3$ の範囲である．
いずれも正である（計算結果~\ref{cr:framing-batch} の一部）．
$x$ を正に大きくすると項数が増え，負では減る（$K10$ の $\rho=\yng(3)$ で，$x=-2$ では約 175 万項，$x=2$ では約 186 万項）．
\end{compresult}
3 箱でも $\bt^{-3x}\hf^{(x)}_\rho$ は 7 根の指数和である（命題~\ref{prop:expsum}）．
しかし，6 根のニュートン形では一様な符号の並びが見つからず，全ての $x$ の証明書は得られていない．

\subsection{予想 B とその動機}\label{sec:conjB}
% [R14]
\begin{tcolorbox}
\begin{conjecture}\label{conj:B}
代数的な iterated torus 結び目は，framing 0 で全ての $\rho$ について正である．
\end{conjecture}
\end{tcolorbox}

KN の物理的な議論は，追加の $U(1)_R$ 対称性を Seifert 多様体の構造から得る．
代数的な iterated torus 結び目の補空間は一般に Seifert 多様体ではないので，この議論は直接は当てはまらない．
予想~\ref{conj:B} の動機は，次の Cherednik--Danilenko~\cite{Cherednik14} の予想にある．
\begin{itemize}
\item 正値性の予想：代数的な結び目と任意の長方形のヤング図形について，DAHA superpolynomial の係数は全て正である（このような結び目を pseudo-algebraic と呼ぶ）．
\item Connection Conjecture (ii)：彩色なしの pseudo-algebraic な結び目では，DAHA superpolynomial は（tilde 正規化した）reduced stable Khovanov--Rozansky 多項式に一致する．
\end{itemize}
これらと本論文の結果の関係は次のとおりである．
\begin{itemize}
\item $\yng(1)$：Connection Conjecture (ii) の下で，DAHA の $\yng(1)$ は KhR の Poincaré 多項式 $P_{\mathrm{KhR}}$ の単項式倍 $mP_{\mathrm{KhR}}$ になる．
(N1) により両者は $\bt=-1$ で同じ reduced HOMFLY-PT 多項式を与えるので，$m=\pm\bt^k$ で $\pm(-1)^k=1$ である．両者とも係数が非負なので $m=\bt^{2l}$ である．
$\bt^{2l}$ は $\sigma$ で不変なので，命題~\ref{prop:one} が適用でき，$\yng(1)$ の正値性が従う．
\item 長方形：Cherednik--Danilenko は DAHA の値の正値性を予想し，ホモロジーとの一致を期待している．予想~\ref{conj:B} の長方形の部分は，これを KN の BPS の段に移したものである．
\item $\yng(2,1)$：Cherednik--Danilenko は長方形でない彩色の正値性を主張していない．実際 $\rCS_{\yng(2,1)}$ は負の係数をもつ．BPS の段での正値性は，その範囲の外の現象である．
\item 非代数的なケーブル：Cherednik--Danilenko は，正値性をもたない非代数的な iterated torus 結び目があることを認めている．本論文の非代数的なケーブルの多くで，DAHA の値の $\yng(1)$ から負の係数が現れる（付録~\ref{app:nonalg}）．
\end{itemize}

\subsection{双曲結び目の 3 箱へ：非長方形の彩色の問題}\label{sec:21}
双曲結び目では，定義に基づく refined $\yng(2,1)$ が文献にもソフトウェアにも無い．
そこで，構造的な条件から $\yng(2,1)$ を決められるかを調べた．
% [R15, R16, R17]
\begin{itemize}
\item 結び目ごと：条件は $t=q$ での 3 本ケーブルから得た $\yng(2,1)$ 彩色 HOMFLY-PT 多項式，ミラー，指数増大性，$U(2)$ への簡約，3 箱の Laurent 性，$G$ 展開である．
これらを課しても $\yng(2,1)$ は一意に決まらない．解の差は「結び目によらない因子 × 自由な多項式」の形をしている．
曖昧さがあること自体は Anokhina ら~\cite{Anokhina12} も指摘している．
\item 弱い判定：正値性を満たす整数の $\yng(2,1)$ は，$\mathbf{4_1}$（KN の候補）と $\mathbf{6_1}$（813 項の解）に存在する．$\mathbf{5_2}$ では有理数の解しか見つからない（整数解の有無は未決）．
これは予想の検証ではなく，予想と矛盾しないことの確認にとどまる．
\item 族として：$t=q$ ではダブルツイスト結び目の $\yng(2,1)$ について Morozov の式~\cite{Morozov16,Morozov19} がケーブルからの値と一致する．これは既知の式の検証である．
refined の長方形は KNTZ~\cite{Kameyama19} の式を実装し，GNSSS の式と三葉結び目の楕円 Hall 代数の値（$\yng(2,2)$ まで）と照合した．しかし refined の $\yng(2,1)$ は，族の構造の仮定を置いても一意に決まらない．
\end{itemize}
双曲結び目の 3 箱には，定義に基づく $\yng(2,1)$ が必要である．その候補は \S\ref{sec:discussion} で述べる．

\section{議論と未解決の問題}\label{sec:discussion}

\paragraph{全ての framing での正値性．}
個々の結び目では，$x$ の指数和とニュートン形によって，全ての $x$ での正値性を有限個の判定に帰着できた（定理~\ref{thm:two-all-x}）．
しかしこれは結び目ごとの判定で，一般の結び目についての定理ではない．
ニュートン形の係数の正値性を，入力の性質（ホモロジーの非負性と $\sigma$ 不変性など）から導けるかは未解決である．
ツイスト結び目で $p$ と $x$ が逆符号の領域と，3 箱の全ての $x$ の証明書も得られていない．
また，Cheng--Sułkowski~\cite{Cheng21} の toric brane の refined open BPS 不変量と，framing をつけた自明な結び目の $\hf_\rho$ の対応は確かめていない．

\paragraph{framing の物理的な導出．}
Mariño--Vafa の符号が Laurent 性に必要であることは示した（命題~\ref{prop:sign}）．
しかし，refined の平坦座標のずれ $V\mapsto-(q/t)^{1/2}V$ を幾何から導くことはしていない．
弦側で $U(1)_R$ の作用が framing によらず自明であることも仮定にとどまる．

\paragraph{双曲結び目の $[2,1]$．}
双曲結び目の 3 箱には，定義に基づく refined $[2,1]$ が必要である．候補は三つある．
(i) refined の Racah 行列による構成（arborescent 計算の refined 化）．
(ii) DAHA・楕円 Hall 代数型の定義の双曲結び目への拡張（種数 2 の DAHA は $\mathfrak{sl}_2$ に限られる~\cite{Arthamonov15}）．
(iii) $[2,1]$ 彩色 HOMFLY-PT ホモロジーの計算と，KN の $\rCS$ への変換規則．
いずれも本論文の範囲を超える．

\paragraph{正規化．}
単項式の正規化の閉じた式（計算結果~\ref{cr:M}）は，refined 衛星公式（仮説~\ref{hyp:sat}）の下で導出した．
DAHA と楕円 Hall 代数は単項式を決めないので，この仮説は物理的な仮説として残る．
トーラス結び目の部分では，Mellit の系の標準の正規化での形を計算で確かめた（定理~\ref{thm:M}）．
$d_\pm$ の表現の元と Bergeron--Garsia--Leven--Xin の $Q_{m,n}$ の対応を具体的に示すことが残っている．
ただし，これは正値性の結果には影響しない．

\paragraph{そのほか．}
pretzel 結び目の異符号の領域の残り（\S\ref{sec:two}）と，2 箱の単峰性の構造的な理由も未解決である．


\appendix
\section{計算法の詳細}\label{app:computation}
\paragraph{評価と補間．}
$\overline{\rCS}_\lambda$ と $\hf_\rho$ は，31 ビットの素数 $p$ を法として格子の点で評価し，$(A,Q,T)$ の各変数について補間して求める．
係数が $p/2$ を超えるときは 2〜3 個の素数で補間し，中国剰余定理（Garner の方法）で復元する．
得られた多項式は，補間に使っていない素数の点で検証する．
評価の点はまとまりごとに保存し，計算を途中から再開できるようにした．
\paragraph{genus 展開．}
補題~\ref{lem:G} と同じく $v=\bq^2\bt$，$w=v+v^{-1}$，$w_0=\bt+\bt^{-1}$ とおくと $G=w-w_0$ である．各行を $w$ の多項式に直し，Horner 法で $G$ の冪に展開する．
\paragraph{計算時間．}
計算は Apple M3 Pro（メモリ 36 GB）と Mac mini M1（メモリ 8 GB）で行った．
最も重いのは $K10$ である．$\rCS_{\yng(3)}$ は 2 つの素数で約 4 時間（14440 秒），$\rCS_{\yng(2,1)}$ は約 2 時間（7387 秒）かかった．
3 箱の $\wh f_\rho$ は 3 つの素数を要し，4050 秒，最大メモリ 17.6 GB であった．
$K7$ と $K8$ の $|\rho|\le3$ は，それぞれ 12275 秒と 22853 秒で，最大メモリは 2 GB 未満であった．
\paragraph{実装の検証．}
実装は，小さな入力で各段を確かめるテストで検証した．
対象は，KN の値の再現，Cherednik--Danilenko の明示例，ブレイドから計算した HOMFLY-PT 多項式，補間と中国剰余定理，genus 展開の二つの実装の一致などである．

\section{正規化の導出}\label{app:normalization}
\paragraph{(A) ケーブル．}
framing 演算子 $\tau:P_\mu\mapsto T_\mu P_\mu$ について，$\tau^jP_\lambda[p_k\mapsto P_{k(m,n)}]\cdot1=P_\lambda[p_k\mapsto P_{k(m,n+jm)}]\cdot1$ が成り立つ（計算結果~\ref{cr:tau}）．
ケーブルの構成は companion について線形である．companion に (N2$'$) を適用すると，パターンの傾きが $(m,-s')$ から $(m,-n)$（Seifert framing）にずれる．
仮説~\ref{hyp:sat} と合わせると，$M_{C_{m,n}K}=M_K^mM_{T_{m,n}}$ と $F=mn$ が従う．
\paragraph{(B) トーラス結び目．}
(I) Gorsky--Neguţ の正規化前の reduced 値の $a^0$ の係数は $(-t/q)^{m-1}C_{m,n}(1/q,t)$ である（$C_{m,n}$ は有理 $q,t$-Catalan 数）．
これは Gorsky--Neguţ の予想であり，Bergeron--Garsia--Leven--Xin の $Q_{m,n}$ を介して Mellit~\cite{Mellit16} の定理から従う．
(II) 最低 $a$ 次数 $\delta$ は Franks--Williams の正のブレイドについての定理から従う．
(III) $t=q$ での符号と $q$ の冪は，$C_{m,n}$ の $q,t$ 対称性から決まる．
(IV) $(q/t)$ の指数は，Mellit~\cite{Mellit17} の Dyck 路の表示と計算結果~\ref{cr:mellit} の式から決まる．
以上から $M_{T_{m,n}}=(-q/t)^{m-1}(a(q/t)^{1/2})^{\delta}$ である．

\section{非代数的なケーブル}\label{app:nonalg}
本論文で扱う非代数的なケーブルは，$T_{2,3}$ の $(2,1)$-，$(2,3)$-，$(2,5)$-，$(2,7)$-，$(2,-1)$-ケーブル（順に $K2,K4,K5a,K5b,K5c$）と，$T_{2,5}$ の $(2,1)$-ケーブル（$K6b$）である．
非代数的なケーブル（$n_2<m_1n_1m_2$）では，DAHA の値はホモロジーの Poincaré 多項式ではない（スペクトル系列の必要条件を満たさない）．
したがって KN の正値性の前提を満たさない．表~\ref{tab:nonalg} に判定を示す．
\begin{table}[h]
\centering
\caption{非代数的なケーブルの判定（framing 0）．各欄は「spin 表示の負の項の数/項数; genus 表示の負の項の数/項数」（全体の符号はそれぞれで良い方）．}\label{tab:nonalg}
\resizebox{\textwidth}{!}{% scripts/paper_summary.py が生成。\checkmark(n) は G 展開の前後とも負の項 0（n は前の項数）。それ以外は 前の負/前の項数; 後の負/後の項数（前後それぞれで全体の符号は良い方）。framing 0。
\begin{tabular}{llcccccc}
\hline
結び目 & 分類 & $\Y{1}$ & $\Y{2}$ & $\Y{1,1}$ & $\Y{3}$ & $\Y{2,1}$ & $\Y{1,1,1}$ \\
\hline
K2 & non & 10/47; 6/48 & 152/1136; 41/1045 & 174/1273; 51/1190 & 967/8525; 253/7852 & 1078/9527; 292/8774 & 1156/9990; 318/9254 \\
K4 & non & 18/42; 20/43 & 286/1109; 195/952 & 330/1255; 224/1090 & 2561/8614; 1749/7420 & 2850/9646; 1989/8325 & 3025/10121; 2128/8795 \\
K5a & non & 8/38; 5/47 & 409/1226; 332/1066 & 460/1377; 376/1210 & 4441/9758; 3631/8278 & 4982/10844; 4066/9235 & 5164/11348; 4263/9733 \\
K5b & non & \checkmark\,(48) & 136/1395; 66/1305 & 149/1537; 78/1452 & 2097/10706; 1060/9612 & 2358/11707; 1181/10553 & 2351/12202; 1228/11040 \\
K5c & non & 5/57; 0/69 & 78/1307; 0/1299 & 77/1443; 0/1433 & 400/9501; 6/9123 & 423/10494; 7/10053 & 376/10864; 0/10444 \\
K6b & non & 30/133; 10/178 & 576/4400; 81/4367 & 618/4677; 91/4656 & 3626/33214; 512/32175 & 3766/35136; 541/34032 & 3878/36067; 560/34980 \\
\hline
\end{tabular}
}
\end{table}
破れ方は三つの型に分かれる．
(I) $T_{2,3}$ の $(2,3)$，$(2,5)$-ケーブル（$K4$，$K5a$）では，色付きでない HOMFLY-PT の段階で，偶奇のそろった refinement による正値性が不可能である．
(II) $(2,1)$-ケーブル（$K2$，$K6b$）と $(2,-1)$-ケーブル（$K5c$）では，HOMFLY-PT は正値性と両立するが，DAHA の $\rCS_\square$ 自体が負の係数をもつ．
(III) $T_{2,3}$ の $(2,7)$-ケーブル（$K5b$）の色付きの場合の破れは framing に依存する．偶数の framing $x$ を調べると，2 箱は $x\le-12$，3 箱は $x\le-14$ で前後とも正になる（$|x|\le24$ の偶数で計算）．
$K5c$ は展開の後だけが（ほぼ）正になる例で，補題~\ref{lem:G} の逆が成り立たないことを示す．

\section{データとコード}\label{app:data}
$\overline{\rCS}_\lambda$，$\hf_\rho$，spin 表示と genus 表示の全ての値と，計算・判定のコード，テストを公開する．


\bibliographystyle{amsplain}
\bibliography{../references}
\end{document}
